% from ega2-1-fr.tex
\oldpage[II]{110}

on a $S'=\operatorname{Spec}(A[c_1,\ldots,c_n])$
(\hyperref[I.6.3.3-fr]{I, 6.3.3}), et $S'$ est donc isomorphe à un
sous-préschéma fermé de $S''=\operatorname{Spec}(A[T_1,\ldots,T_n])$ ($T_i$
indéterminées). Dans le diagramme commutatif
\[
\xymatrix{
X\times_SS'\ar[r]^{1_X\times j}\ar[d]_{(f_0)_{(S')}} &
X\times_SS''\ar[d]^{(f_0)_{(S'')}}\\
S'\ar[r]_j & S''
}
\]
$j$ et $1_X\times j$ sont des immersions fermées
(\hyperref[I.4.3.1-fr]{I, 4.3.1}) et $(f_0)_{(S'')}$ est fermé par hypothèse~;
il en est donc de même de $(f_0)_{(S')}$. 

\section{Morphismes entiers et morphismes finis}
\label{section:II.6-fr}

\subsection{Préschémas entiers sur un autre.}
\label{subsection:II.6.1-fr}

\begin{definition}[6.1.1]
\label{II.6.1.1-fr}
Soit $X$ un $S$-préschéma, $f:X\to S$ son morphisme structural. On dit que
$X$ est \emph{entier au-dessus de $S$}, ou que $f$ est un \emph{morphisme
entier}, s'il existe un recouvrement $(S_\alpha)$ de $S$ par des ouverts
affines tels que, pour tout $\alpha$, le préschéma induit $f^{-1}(S_\alpha)$
soit un schéma affine, dont l'anneau $B_\alpha$ est une algèbre entière sur
l'anneau $A_\alpha$ de $S_\alpha$. On dit que $X$ est \emph{fini au-dessus de
$S$}, ou que $f$ est un \emph{morphisme fini} si $X$ est entier sur $S$ et de
type fini sur $S$.
\end{definition}

Lorsque $S$ est affine d'anneau $A$, on dira aussi «~entier (resp. fini) sur
$A$~» au lieu de «~entier (resp. fini) sur $S$~».

\begin{env}[6.1.2]
\label{II.6.1.2-fr}
Il est clair que si $X$ est entier au-dessus de $S$, il est affine au-dessus
de $S$. Pour qu'un préschéma $X$ affine au-dessus de $S$ soit entier (resp.
fini) au-dessus de $S$, il faut et il suffit que la $\mathcal{O}_S$-Algèbre
quasi-cohérente associée $\mathcal{A}(X)$ soit telle qu'il existe un
recouvrement $(S_\alpha)$ de $S$ par des ouverts affines ayant la propriété
que pour tout $\alpha$, $\Gamma(S_\alpha,\mathcal{A}(X))$ est une algèbre
entière (resp. entière et de type fini) sur
$\Gamma(S_\alpha,\mathcal{O}_S)$. Une $\mathcal{O}_S$-Algèbre quasi-cohérente
ayant cette propriété est dite \emph{entière} (resp. \emph{finie}) sur
$\mathcal{O}_S$. Se donner un préschéma entier (resp. fini) au-dessus de $S$
revient donc (\hyperref[II.1.3.1-fr]{1.3.1}) à se donner une
$\mathcal{O}_S$-Algèbre quasi-cohérente entière (resp. finie) sur
$\mathcal{O}_S$. On notera qu'une $\mathcal{O}_S$-Algèbre quasi-cohérente
$\mathcal{B}$ est finie si et seulement si c'est un $\mathcal{O}_S$-Module de
type fini (\hyperref[I.1.3.9-fr]{I, 1.3.9})~; il revient au même de dire que
$\mathcal{B}$ est une $\mathcal{O}_S$-Algèbre entière et de type fini, car une
algèbre entière et de type fini sur un anneau $A$ est un $A$-module de type
fini.
\end{env}

\begin{proposition}[6.1.3]
\label{II.6.1.3-fr}
Soit $S$ un préschéma localement noethérien. Pour qu'un préschéma $X$ affine
au-dessus de $S$ soit fini au-dessus de $S$, il faut et il suffit que la
$\mathcal{O}_S$-Algèbre $\mathcal{A}(X)$ soit cohérente.
\end{proposition}

Compte tenu de la remarque précédente, cela revient à remarquer que si $S$ est
localement noethérien, les $\mathcal{O}_S$-Modules quasi-cohérents de type fini
ne sont autres que les $\mathcal{O}_S$-Modules cohérents
(\hyperref[I.1.5.1-fr]{I, 1.5.1}).

\begin{proposition}[6.1.4]
\label{II.6.1.4-fr}
Soit $X$ un préschéma entier (resp. fini) au-dessus de $S$, $f:X\to S$ le
morphisme structural. Alors, pour tout ouvert affine $U\subset S$, d'anneau
$A$, $f^{-1}(U)$ est un schéma affine dont l'anneau $B$ est une algèbre
entière (resp. finie) sur $A$.
\end{proposition}

